\subsection{投射模的秩}

定义 1. 设 P 是有限生成投射 A-模。对 A 的每个素理想 p，自由 \(A_{p}\) -模 \(P_{p}\) 的秩称为 P 在 p 处的秩，记作 \(\operatorname{rg}_{\mathfrak{p}}(\mathbf{P})\)。

由定理 1，整值函数 \(\mathfrak{p} \mapsto \mathrm{rg}_{\mathfrak{p}}(\mathrm{P})\) 在 \(\mathbf{X} = \operatorname{Spec}(\mathbf{A})\) 上是局部常数的；因此当 \(\mathbf{X}\) 连通时，特别地当环 \(\mathbf{A}\) 是整环时，它是常数（§ 4, 第 3 节, 命题 15 的推论 2）。

定义 2. 设 \(n\) 是一个整数 \(\geqslant 0\) 。称投射 A-模 \(P\) 是秩 \(n\) 的，如果它有限生成，并且关系式 \(\operatorname{rg}_{\mathfrak{p}}(P) = n\) （其中 \(\mathfrak{p}\) 取遍 \(A\) 的素理想）成立。

显然，每个有限生成的自由 A-模 L 依定义 2 的意义都是秩为 n 的模，这里 n 等于《代数》第 II 章 §7 第 2 节中定义的 L 的维数（或秩）。

秩为 0 的投射模是零模（§ 3, 第 3 节, 定理 1 的推论 2）。若 A 不化为 0，且投射 A-模 P 的秩为 n，则整数 n 被唯一确定；这时把它记作 rg(P)。

定理 2. 设 P 是一个 A-模，n 是整数 \(\geqslant 0\) 。下列性质等价：

\begin{enumerate}
\def\labelenumi{(\alph{enumi})}
\item
  P 是秩为 n 的投射模。
\item
  \(\mathbf{P}\) 有限生成，且对每个极大理想 \(\mathfrak{m}\) （在 \(\mathbf{A}\) 中），\(\mathbf{A}_{\mathfrak{m}}\) -模 \(\mathbf{P}_{\mathfrak{m}}\) 是秩为 \(n\) 的自由模。
\item
  \(\mathbf{P}\) 有限生成，且对每个素理想 \(\mathfrak{p}\) （在 \(\mathbf{A}\) 中），\(\mathbf{A}_{\mathfrak{p}}\) -模 \(\mathbf{P}_{\mathfrak{p}}\) 是秩为 \(n\) 的自由模。
\item
  对 A 的每个极大理想 m，存在 \(f \in A - m\) ，使得 \(A_{f}\) -模 \(P_{f}\) 是秩为 n 的自由模。
\end{enumerate}

由定义 2 与定理 1，(a) 与 (c) 等价；(b) 蕴含 (c)，因为对 A 的每个素理想 p，存在包含 p 的极大理想 m，并且记 \(p' = p_{m}\) 时，\(P_{p}\) 同构于 \((\mathrm{P}_{\mathfrak{m}})_{\mathfrak{p'}}\)（§ 2, 第 5 节, 命题 11）；若 \(P_{m}\) 是秩为 n 的自由模，则 \(P_{p}\) 也是。性质 (c) 依据定理 1 以及如下事实蕴含 (d)：若 \(f \in A - m\) 且 \(m' = m_{f}\) ，则 \(P_{m}\) 同构于 \((\mathrm{P}_{f})_{\mathfrak{m'}}\) ，因而 \(P_{f}\) 与 \(P_{m}\) 的秩相等。最后，上面的论证与定理 1 表明 (d) 蕴含 (b)。

注. 若 A 是整环，则如上文所见，投射 A-模 具有完全确定的秩（依定义 2 的意义）；而且这个秩与《代数》第 II 章 §7 第 2 节中定义的秩一致；只需对 p = (0) 应用定理 2 (c) 即可。

设 E 与 F 是两个有限生成的投射 A-模。我们知道（《代数》第 II 章 §§ 2 与 3），\(\mathbf{E} \times \mathbf{F}\) 、\(\mathbf{E} \otimes_{\mathbf{A}} \mathbf{F}\) 、\(\operatorname{Hom}_{\mathbf{A}}(\mathbf{E}, \mathbf{F})\) 以及 E 的对偶 \(\mathbf{E}^*\)

（即 E 的对偶模）都是投射的且有限生成的；外幂 \(\wedge^k\) E 对每个整数 \(k > 0\) 亦然（《代数》第 III 章）。另外，由定义 1 以及 §2 第 7 节的命题 18、19 和第 8 节立即可得，对 A 的每个素理想 \(\mathfrak{p}\) ：

\[
\mathrm{rg} _ {\mathfrak {p}} (\mathrm{E} \times \mathrm{F}) = \mathrm{rg} _ {\mathfrak {p}} (\mathrm{E}) + \mathrm{rg} _ {\mathfrak {p}} (\mathrm{F})\tag{1}
\]

\[
\operatorname{rg} _ {\mathfrak {p}} (\mathrm{E} \otimes_ {\mathrm{A}} \mathrm{F}) = \operatorname{rg} _ {\mathfrak {p}} (\mathrm{E}). \operatorname{rg} _ {\mathfrak {p}} (\mathrm{F})\tag{2}
\]

\[
\operatorname{rg} _ {\mathfrak {p}} (\operatorname{Hom} _ {\mathbf {A}} (\mathrm{E}, \mathrm{F})) = \operatorname{rg} _ {\mathfrak {p}} (\mathrm{E}). \operatorname{rg} _ {\mathfrak {p}} (\mathrm{F})\tag{3}
\]

\[
\operatorname{rg} _ {\mathfrak {p}} (\mathrm{E} ^ {*}) = \operatorname{rg} _ {\mathfrak {p}} (\mathrm{E})\tag{4}
\]

\[
\operatorname{rg} _ {\mathfrak {p}} (\overset {k} {\wedge} \mathrm{E}) = \binom{\operatorname{rg} _ {\mathfrak {p}} (\mathrm{E})}{k}.\tag{5}
\]

若 E 与 F 的秩都有定义，则 \(E \times F\) 、\(E \otimes_{A} F\) 、\(\operatorname{Hom}_{A}(E, F)\) 、\(E^{*}\) 与 \(\Lambda^{k} E\) 的秩也都 有定义，而且上面的等式在略去指标 p 后仍然成立。此外：

推论. 有限生成投射 A-模 P 是秩为 n 的模，其充分必要条件是 \(\wedge^{n}P\) 的秩为 1。

命题 4. 设 B 是交换 A-代数，P 是秩为 n 的投射 A-模。则 B-模 \(\mathbf{P}_{(\mathbf{B})} = \mathbf{B} \otimes_{\mathbf{A}} \mathbf{P}\) 是秩为 n 的投射模。

我们知道 \(P_{(B)}\) 是投射的且有限生成的（《代数》第 II 章 §5 第 1 节命题 4 的推论）。若 q 是 B 的素理想，p 是它在 A 中的逆像，则

\[
\left(\mathrm{P} _ {(\mathrm{B})}\right) _ {\mathrm{q}} = \left(\mathrm{P} \otimes_ {\mathrm{A}} \mathrm{B}\right) \otimes_ {\mathrm{B}} \mathrm{B} _ {\mathrm{q}} = \mathrm{P} \otimes_ {\mathrm{A}} \mathrm{B} _ {\mathrm{q}} = \left(\mathrm{P} \otimes_ {\mathrm{A}} \mathrm{A} _ {\mathrm{p}}\right) \otimes_ {\mathrm{A}} \mathrm{B} _ {\mathrm{q}}
\]

而且，由假设，\(\mathbf{P} \otimes_{\mathbf{A}} \mathbf{A}_{\mathfrak{p}}\) 是自由 \(\mathbf{A}_{\mathfrak{p}}\) -模且秩为 \(n\) ，故 \((\mathrm{P}_{(\mathbf{B})})_{\mathfrak{q}}\) 是自由 \(\mathbf{B}_{\mathfrak{q}}\) -模且秩为 \(n\) 。

命题 5. 设 A 是半局部环，P 是有限生成的投射 A-模。若 P 的秩有定义，则 P 是自由 A-模。

首先假设 A 同构于域 \(K_{i}\) ( \(1 \leqslant i \leqslant n\) ) 的积。这时 \(K_{i}\) 等同于 A 的极小理想（《代数》第 VIII 章 §3 第 1 节），而且对每个 i，和 \(p_{i}\) （它是诸 \(K_{j}\) （ \(j \neq i\) ）之和）是 A 的极大理想，\(p_{i}\) ( \(1 \leqslant i \leqslant n\) ) 是 A 仅有的素理想。因此每个有限生成 A-模 P 是其同型分量 \(P_{i}\) ( \(1 \leqslant i \leqslant n\) ) 的直和，其中 \(P_{i}\) 同构于有限多个（\(r_{i}\) 个）同构于 \(K_{i}\) 的 A-模的直和（《代数》第 VIII 章 §5 第 1 节命题 1 与第 3 节命题 11）；环 \(A_{p_{i}}\) 等同于 \(K_{i}\) ，并零化诸 \(P_{j}\) （ \(j \neq i\) ） ，由此 \(r_{i} = \operatorname{rg}_{\mathfrak{p}_{i}}(\mathbf{P})\) ；若所有 \(r_{i}\) 都等于同一个数 r，则 P 同构于 \(A^{r}\) ，命题在这一情形得证。在一般情形，设 R 是 A 的 Jacobson 根，B = A/R；由于 B 是域的积，由命题 4 之前的注可知投射 B-模 \(P_{(B)}\) 是自由的。又 P 是平坦 A-模，于是命题由 §3 第 2 节命题 5 得出。

